\chapter{The model and the classes}
\label{chap:model}

The argument is stated once for an abstract model and once for Turing machines. The abstract
model only names the classes and the closure properties that the argument uses; the
Turing-machine instance takes them from cslib and cslib-frontier, with their proofs.

\section{The abstract model}

\begin{definition}[Abstract complexity model]
  \label{def:model}
  \lean{Sagitta.Arrow.Model}
  \leanok
  A \emph{model} consists of a type of decision problems and a type of search problems; predicates
  $\cP \subseteq \cNP$ on decision problems and $\cFP$, $\cFNP$ on search problems; a reflexive and
  transitive relation $A \red B$ of reductions; and two maps $\dom$ and $\pre$ from search problems
  to decision problems (the domain, ``does the instance have a solution?'', and the prefix problem,
  ``does some solution of $a$ begin with $u$?''). They satisfy:
  \begin{enumerate}
    \item if $A \red B$ and $B \in \cP$ then $A \in \cP$ (an efficient procedure may call an
      efficient subroutine);
    \item if $R \in \cFNP$ then $\dom R \in \cNP$ and $\pre R \in \cNP$; if $R \in \cFP$ then
      $\dom R \in \cP$;
    \item every $A \in \cNP$ is $\dom R$ for some $R \in \cFNP$;
    \item if $R \in \cFNP$ and $\pre R \in \cP$ then $R \in \cFP$ (search through prefixes).
  \end{enumerate}
  \own{} The closure properties are standard facts about the classes of the literature: see
  \cite{AroraBarak} (Chapter~2) and \cite{Goldreich} (Chapter~2, search and decision).
\end{definition}

\begin{definition}[$\cP = \cNP$, easy search, completeness, intermediate problems]
  \label{def:model-classes}
  \lean{Sagitta.Arrow.Model.PeqNP, Sagitta.Arrow.Model.SearchEasy, Sagitta.Arrow.Model.NPComplete, Sagitta.Arrow.Model.Intermediate}
  \leanok
  \uses{def:model}
  In a model $M$:
  $\cP = \cNP$ means that every $A \in \cNP$ is in $\cP$;
  \emph{search is easy} means that every $R \in \cFNP$ is in $\cFP$;
  $B$ is \emph{$\cNP$-complete} when $B \in \cNP$ and $A \red B$ for every $A \in \cNP$;
  $A$ is \emph{$\cNP$-intermediate} when $A \in \cNP$, $A \notin \cP$ and $A$ is not
  $\cNP$-complete.
  \own{}
\end{definition}

\section{Imported from cslib}
\label{sec:cslib}

The entries of this section are \textbf{not this project's work}. They are definitions and
theorems of \href{https://github.com/leanprover/cslib}{cslib} (commit \texttt{98f9a62}) and
\href{https://github.com/cslib-community/cslib-frontier}{cslib-frontier} (commit \texttt{99f0281}),
copied into this repository without modification, under their original paths and module names,
with their Apache-2.0 licenses and copyright headers (see
\href{https://github.com/julatec/sagitta/blob/main/VENDORED.md}{\texttt{VENDORED.md}}). Each entry
says where the copied file is. They are listed here so that the dependency graph shows on which
part of cslib each result of this project rests.

All problems are over instances $a$ of a type $\alpha$ with an injective binary encoding
$\enc : \alpha \hookrightarrow \{0,1\}^*$.

\subsection{The machines}

\begin{definition}[Multi-tape Turing machines]
  \label{def:cslib-tm}
  \lean{Turing.MultiTapeTM, Turing.Cfg, Turing.Action, Turing.MultiTapeTM.runFrom}
  \leanok
  A deterministic multi-tape machine with $k$ work tapes has a transition that, from a state, the
  symbol under the input head and the symbols under the work heads, gives an action: move the input
  head, write and move on each work tape, append at most one symbol to the output, and go to a new
  state or halt. A configuration records the state, the head positions, the work tapes and the
  output; $\mathrm{runFrom}(c, t)$ is the configuration after $t$ steps.
  \cslibfile{Cslib/Computability/Machines/Turing/MultiTape/Deterministic.lean}
  (configurations and actions in
  \cslibfile{Cslib/Computability/Machines/Turing/MultiTape/Configuration.lean}).
  \notours{cslib}
\end{definition}

\begin{definition}[Time and space of a run]
  \label{def:cslib-time-space}
  \lean{Turing.MultiTapeNTM.RunsInTime, Turing.MultiTapeNTM.RunsInSpace, Turing.MultiTapeTM.ComputesFunInTimeAndSpace}
  \leanok
  \uses{def:cslib-tm}
  A machine runs in time $t$ on an input when every run halts within $t$ steps, and in space $s$
  when it stays within $s$ cells; it computes $f$ in time $t$ and space $s$ (for encodings
  $\enc_{\mathrm{in}}, \enc_{\mathrm{out}}$) when on $\enc_{\mathrm{in}}(a)$ it halts with output
  $\enc_{\mathrm{out}}(f(a))$ within those bounds.
  \cslibfile{Cslib/Computability/Machines/Turing/MultiTape/Nondeterministic.lean},
  \cslibfile{Cslib/Computability/Machines/Turing/MultiTape/Space.lean} and
  \cslibfile{Cslib/Computability/Machines/Turing/MultiTape/Deterministic.lean}.
  \notours{cslib}
\end{definition}

\begin{definition}[Bounded configurations]
  \label{def:cslib-configbound}
  \lean{Turing.Cfg.core, Turing.Storage, Turing.storageBound}
  \leanok
  \uses{def:cslib-tm}
  The core of a configuration forgets the output; a configuration whose work tapes fit in windows of
  size $s$ is one of at most $\mathrm{storageBound}$ storages, a number exponential in $s$.
  \cslibfile{Cslib/Computability/Machines/Turing/MultiTape/ConfigBound.lean}.
  \notours{cslib}
\end{definition}

\begin{definition}[Plumbing: sequential composition, virtual input and output, rewinding]
  \label{def:cslib-plumbing}
  \lean{Turing.MultiTapeTM.seq, Turing.MultiTapeTM.inputFromTape, Turing.MultiTapeTM.outputToTape, Turing.MultiTapeTM.rewindWork}
  \leanok
  \uses{def:cslib-tm}
  $\mathrm{seq}$ runs one machine and then another on the same tapes; $\mathrm{inputFromTape}$ makes
  a machine read its input from a work tape; $\mathrm{outputToTape}$ writes its output on a fresh
  work tape; $\mathrm{rewindWork}$ brings a work head back to the start of its word.
  \cslibfile{Cslib/Computability/Machines/Turing/MultiTape/Plumbing/Sequential.lean},
  \cslibfile{Cslib/Computability/Machines/Turing/MultiTape/Plumbing/InputFromTape.lean},
  \cslibfile{Cslib/Computability/Machines/Turing/MultiTape/Plumbing/OutputToTape.lean},
  \cslibfile{Cslib/Computability/Machines/Turing/MultiTape/Plumbing/RewindWork.lean}.
  \notours{cslib}
\end{definition}

\subsection{The classes}

\begin{definition}[Polynomial-time computable function]
  \label{def:polytime}
  \lean{Turing.MultiTapeTM.ComputableInPolyTime}
  \leanok
  \uses{def:cslib-time-space}
  $f : \alpha \to \beta$ is computable in polynomial time when some machine with binary alphabet and
  finitely many states computes it within time and space bounded by polynomially bounded functions
  of the length of the encoded input.
  \cslibfile{CslibFrontier/Computability/Machines/Turing/MultiTape/PolyTime.lean}.
  \notours{cslib-frontier}
\end{definition}

\begin{definition}[The class $\cP$]
  \label{def:P}
  \lean{Cslib.Complexity.P}
  \leanok
  \uses{def:polytime}
  $A \in \cP$ when some Boolean function computable in polynomial time decides membership in $A$.
  \cslibfile{CslibFrontier/Computability/Complexity/Classes/P.lean}.
  \notours{cslib-frontier}
\end{definition}

\begin{definition}[The class $\cNP$, by verifiers]
  \label{def:NP}
  \lean{Cslib.Complexity.IsVerifier, Cslib.Complexity.NP}
  \leanok
  \uses{def:polytime}
  $V$ verifies $A$ with polynomial $p$ when
  $a \in A \iff \exists w,\ |w| = p(|\enc(a)|) \wedge V(a, w) = \mathrm{true}$; $A \in \cNP$ when it
  has a verifier computable in polynomial time.
  \cslibfile{CslibFrontier/Computability/Complexity/Classes/NP.lean}.
  \notours{cslib-frontier}
\end{definition}

\begin{definition}[Polynomial-time reductions and $\cNP$-completeness]
  \label{def:reduction}
  \lean{Cslib.Complexity.PolyTimeReducible, Cslib.Complexity.PolyTimeReducible.trans, Cslib.Complexity.NPComplete, Cslib.Complexity.P_subset_NP}
  \leanok
  \uses{def:P, def:NP}
  $A \red B$ when some $f$ computable in polynomial time satisfies $a \in A \iff f(a) \in B$;
  reductions compose. $B$ is $\cNP$-complete when $B \in \cNP$ and every $\cNP$ problem reduces to
  it. $\cP \subseteq \cNP$.
  \cslibfile{CslibFrontier/Computability/Complexity/Reduction.lean}.
  \notours{cslib-frontier}
\end{definition}

\begin{definition}[Search problems: $\cFNP$ and $\cFP$]
  \label{def:FNP}
  \lean{Cslib.Complexity.FNP, Cslib.Complexity.FPSearch}
  \leanok
  \uses{def:polytime}
  A search problem is a relation $R(a, w)$. $R \in \cFNP$ when its solutions are polynomially short
  and it is decided in polynomial time; $R \in \cFP$ when moreover a function computable in
  polynomial time returns a solution whenever there is one.
  \cslibfile{CslibFrontier/Computability/Complexity/Classes/TFNP.lean}.
  \notours{cslib-frontier}
\end{definition}

\begin{theorem}[Search reduces to decision]
  \label{thm:cslib-search-decision}
  \lean{Cslib.Complexity.FNP_subset_FPSearch_iff}
  \leanok
  \uses{def:FNP, def:P, def:NP}
  $\cFNP \subseteq \cFP$ for every encoding if and only if $\cNP \subseteq \cP$ for every encoding.
  \cslibfile{CslibFrontier/Computability/Complexity/Classes/TFNP/SearchToDecision.lean}.
  \notours{cslib-frontier}
\end{theorem}

\begin{proof}
  \leanok
  In cslib-frontier.
\end{proof}

\subsection{The problems}

\begin{definition}[$\SAT$ and $\TSAT$]
  \label{def:threesat}
  \lean{Cslib.Complexity.SAT, Cslib.Complexity.encCNF, Cslib.Complexity.ThreeSAT}
  \leanok
  $\SAT$ is the set of satisfiable CNF formulas over $\mathbb{N}$, with a binary encoding of the
  clause list; $\TSAT$ the satisfiable ones with at most three literals per clause.
  \cslibfile{CslibFrontier/Computability/Complexity/Problems/SAT.lean},
  \cslibfile{CslibFrontier/Computability/Complexity/Problems/ThreeSAT.lean}.
  \notours{cslib-frontier}
\end{definition}

\begin{theorem}[Cook--Levin]
  \label{thm:cook-levin}
  \lean{Cslib.Complexity.CookLevin.SAT_npComplete, Cslib.Complexity.ThreeSAT_npComplete, Cslib.Complexity.polyTimeReducible_SAT_ThreeSAT}
  \leanok
  \uses{def:threesat, def:reduction}
  $\SAT$ and $\TSAT$ are $\cNP$-complete, and $\SAT \red \TSAT$.
  \cslibfile{CslibFrontier/Computability/Complexity/CookLevin/Completeness.lean},
  \cslibfile{CslibFrontier/Computability/Complexity/Problems/ThreeSAT.lean}.
  \notours{cslib-frontier}
\end{theorem}

\begin{proof}
  \leanok
  In cslib-frontier.
\end{proof}

\begin{theorem}[$\MAXCUT$ is $\cNP$-complete]
  \label{thm:maxcut}
  \lean{Cslib.Complexity.MaxCut, Cslib.Complexity.MaxCut_npComplete}
  \leanok
  \uses{def:reduction}
  The weighted maximum-cut problem (is there a cut of weight at least $k$?) is $\cNP$-complete.
  \cslibfile{CslibFrontier/Computability/Complexity/Problems/MaxCut.lean}.
  \notours{cslib-frontier}
\end{theorem}

\begin{proof}
  \leanok
  In cslib-frontier.
\end{proof}

\begin{theorem}[Ladner]
  \label{thm:ladner}
  \lean{Cslib.Complexity.Ladner.exists_NP_not_P_not_npComplete}
  \leanok
  \uses{def:P, def:NP, def:reduction}
  If $\cP \neq \cNP$ for some encoding, some problem over binary words is in $\cNP$, not in $\cP$
  and not $\cNP$-complete.
  \cslibfile{CslibFrontier/Computability/Complexity/Classes/Ladner.lean}.
  \notours{cslib-frontier}
\end{theorem}

\begin{proof}
  \leanok
  In cslib-frontier.
\end{proof}

\section{Turing machines form a model}

\begin{definition}[The Turing-machine model]
  \label{def:tmModel}
  \lean{Sagitta.Arrow.TM.tmModel, Sagitta.Arrow.TM.Problem, Sagitta.Arrow.TM.SearchProblem}
  \leanok
  \uses{def:model, def:P, def:NP, def:reduction, def:FNP, thm:cslib-search-decision}
  Decision problems are triples $(\alpha, \enc, A)$, search problems triples $(\alpha, \enc, R)$,
  reductions are polynomial-time many-one reductions, and the classes are those of cslib-frontier.
  Every closure property of Definition~\ref{def:model} is a theorem of cslib-frontier, used with its
  proof: composition of reductions, closure of $\cP$ under reductions, the domain of an $\cFNP$
  relation in $\cNP$, every $\cNP$ problem the domain of its verifier relation, and search through
  prefixes.
\end{definition}

\begin{lemma}[Completeness in the two vocabularies]
  \label{lem:npComplete-iff}
  \lean{Sagitta.Arrow.TM.npComplete_iff}
  \leanok
  \uses{def:tmModel, def:model-classes}
  For a problem $(\beta, \enc, B)$, $\cNP$-completeness in the abstract model coincides with
  cslib's $\cNP$-completeness.
\end{lemma}

\begin{proof}
  \leanok
  Unfold both definitions; an $\cNP$ problem of the model is an $\cNP$ problem over some encoded type.
\end{proof}
